\clearpage
\appendix
\section{What Contrasts Can Identify in a Small Submodel}

It is helpful to distinguish a likelihood that becomes more informative from
one that identifies every latent mechanism. Consider a static, noise-free
submodel of one learner, one form, and one response mode. Let \(b\) be the
base-recall probability, \(p\) the pattern-execution probability, and \(r\)
the whole-form memory probability. Let the three tasks have exact population
response probabilities
\begin{equation}
 p_B=b,\qquad
 p_C=r+(1-r)p,\qquad
 p_F=r+(1-r)bp .
\end{equation}
Here \(B\) is base recall, \(C\) supplies the base and asks for the derived
form, and \(F\) supplies only the lexical action and form category.
If \(b\ne1\) and \(r\ne1\), the stipulated model permits
\begin{equation}
 r=\frac{p_F-bp_C}{1-b},\qquad
 p=\frac{p_C-r}{1-r}.
\end{equation}
This algebra explains why changing the tested ingredients can help.
It also exposes boundaries. When \(b\) is close to one, separating the memory
route from the pattern route is ill-conditioned. If all three response
probabilities are noisy estimates, the algebra need not yield admissible
probabilities without an appropriate likelihood. Different response modes
cannot be substituted for each other in these equations. A recognition
success is not an observation of \(p_B\) for production.

The full implemented model adds contextual, word, language, temporal, and
hierarchical variation; its generic diagnostic flags do not claim the above
three-equation identification result applies automatically. A different base
with the same bounded pattern supplies additional contrast while retaining
its own word memory and lexical knowledge.

\section{Why Inactive Coefficients Can Be Collapsed}

Partition the Gaussian effect vector into coefficients \(\theta_A\) that
appear in the accepted likelihood and coefficients \(\theta_U\) that do not.
Conditional on the variance groups, the declared prior factors over
coefficients. Therefore
\begin{equation}
\begin{split}
 p(\theta_A,\tau^2\mid D)
 &\propto p(D\mid\theta_A)\,
          p(\theta_A\mid\tau^2)\,p(\tau^2)\\
 &\quad\cdot\int p(\theta_U\mid\tau^2)\,d\theta_U\\
 &=p(D\mid\theta_A)\,
          p(\theta_A\mid\tau^2)\,p(\tau^2).
\end{split}
\end{equation}
The equality holds up to normalization. Sampling this marginal posterior
avoids adding many unobserved catalogue
effects to the variance-Gibbs transition. For a retained
\((\theta_A,\tau^2)\), drawing \(\theta_U\) from its conditional Gaussian
prior reconstructs a draw from the full posterior. The same reconstruction
is used across all task predictions for that draw; unrelated marginal
prediction arrays are never independently shuffled before computing EVSI.

The active set is defined by the accepted likelihood, not merely by a task
being present in the catalogue. Variance groups with no active coefficients
remain prior-only. Such labels distinguish a learned shared ability from a
group variance that has evidence elsewhere, and from a wholly untouched
prior. The export preserves all three cases.

\section{Reproduction and Artifact Layout}

The package retains the original graph, reader, account store, and backend.
The new pure modules separate tasks, audited input, posterior inference,
and decisions. The live UI writes through the existing content-event API.
No additional runtime package or external model provider is needed for
inference. Python and browser development dependencies are used only for
existing application hosting, testing, or paper generation.

\begin{table*}[h]
\caption{Implementation entry points and the contract each owns.}
\centering\small
\begin{tabular}{@{}p{.43\textwidth}p{.51\textwidth}@{}}
\toprule
Artifact & Responsibility \\
\midrule
\artifact{public/nabu-learner-tasks.js} &
Versioned task conditions, answer rules, component map, exposure identities,
source integrity gates. \\
\artifact{public/nabu-learner-observations.js} &
Actor-scoped evidence admission, conditioned tasks, exposure reconciliation,
bounded audit witnesses. \\
\artifact{public/nabu-learner-model.js} &
Hierarchical likelihood, collapsed sampling, prediction, hyperparameters,
diagnostics, explicit cohort API. \\
\artifact{public/nabu-learner-policy.js} &
Validated support effects, posterior loss, binary EVSI, budget and work bounds. \\
\artifact{public/nabu-learner-ui.js} &
Goal selection, live task delivery, idempotent writes, cancellation,
posterior refresh and export. \\
\artifact{tools/evaluate_learner.mjs} &
Independent reference calculation and seeded synthetic decision/prediction
report. \\
\artifact{docs/LEARNER_VERIFICATION.md} &
Recorded test outcomes and concrete limits of the execution environment. \\
\bottomrule
\end{tabular}
\end{table*}

Run the implementation and verification commands from the package root:
\begin{verbatim}
python start_local.py
npm test
python -m pytest -q
npm run test:learner
npm run evaluate:learner
python tools/build_standalone.py
\end{verbatim}
The DOM suites require the pinned jsdom development dependency, installed by
\texttt{npm install}; an existing installation can be selected with
\texttt{NABU\_TEST\_JSDOM}. The corresponding commands are
\texttt{npm run test:learner-ui} and the existing evidence, intent, and
next-step UI scripts. This test dependency is not shipped in the application's
runtime bundle.

Generate figures and the current paper from saved results:
{\footnotesize
\begin{verbatim}
python tools/build_learner_paper_figures.py
python tools/build_paper.py \
  --paper nabu-learner
\end{verbatim}
}
The figure manifest records the numerical report hash. The shipping HTML is
rebuilt with exact script and style hashes and a matching service-worker build
identity. The delivery manifest records file hashes for the release package
and paper. Earlier papers and baseline evaluation tools are preserved as
historical artifacts, not presented as results of this new learner model.

\section{Privacy and Practical Boundaries}

The current actor and selected team context delimit usable events. Separate
people are not pooled merely because their accounts belong to one team.
The cohort API is an explicit analytical facility whose caller must supply
authorized data and scope. A downloaded model record contains that actor's
scoped evidence and inferred posterior; it is an intentional user action,
not automatic telemetry.

The existing API accepts authenticated client observations. It does not
proctor performance or verify that a person answered unaided outside the
interface. Exact references and exposure controls make a record interpretable;
they cannot prevent all unrecorded help. The authored 30-minute exposure
window is a conservative operational rule, not an estimated forgetting time
or a claim that responses become statistically independent at its boundary.

Hint events are committed before displaying the cue. If the action is
cancelled while the save is pending, that event conservatively records
possibly supplied support even though the cue may never appear. This can
exclude a later response for the session. The present schema also lacks
the exact selected cue identity, so replay uses the union of maintained
possible cues for that task. Neither case is treated as proof of unaided
success.
